% TOP_PROBLEM: {"id":30007758,"problem_number":"LOCAL-30007758","title":"Tax incidence anomalies","category_id":21,"category":{"id":21,"name":"economics","display_name":"Economics","description":"Open economic theory statements, published research agendas, and proposed scoped specifications; question provenance and historical priority are qualified per record.","slug":"economics","order_index":21,"created_at":"2026-10-08T00:00:00Z"},"status":"provenance_pending","external_url":null,"legacy_ids":[],"published":false,"statement":"In the explicitly specified setting: Why do tax increases and cuts pass through differently to prices and wages, and when do those asymmetries persist? The mathematical statement above fixes the target, assumptions and scope of this catalogue formulation.","economics":{"original_id":247,"field":"Taxation and social insurance","kind":"identification","formalization_basis":"proposed_specification","source_status":"not_verified","priority_status":"not_established","priority_note":"An original explicit open-question statement was not verified; the cited supporting paper must not be treated as its first listing.","earliest_verified_reference":null,"current_reference":{"authors":["Youssef Benzarti"],"title":"Tax Incidence Anomalies","year":2025,"url":"https://www.annualreviews.org/content/journals/10.1146/annurev-economics-081324-085805","doi":"10.1146/annurev-economics-081324-085805","locator":"Annual Review of Economics 17, 615–634; abstract","evidence_summary":"Lists evidence challenging statutory-incidence invariance, elasticity-only incidence and symmetry. The accessed abstract does not explicitly declare the catalogue mechanism problem open. A 2024 working-paper version exists."},"formalization_note":"Catalogue-authored scoped formalization. Population, instruments, welfare objective and identification restrictions must be instantiated before claiming a resolution; source authors are not credited with these added assumptions."},"statusQualification":"Provenance pending: an explicit published statement of this exact open question was not verified. This is a catalogue-authored candidate specification, not a certified open conjecture. Catalogue-authored scoped formalization. Population, instruments, welfare objective and identification restrictions must be instantiated before claiming a resolution; source authors are not credited with these added assumptions.","statusReviewedAt":"2026-10-08"}
% FIRST_POSTED: 2026-10-08
% ENTRY_KIND: statement_only
% !TeX program = lualatex
\documentclass[11pt]{article}
\usepackage{fontspec}
\usepackage[a4paper,margin=25mm]{geometry}
\usepackage{amsmath,amssymb,mathtools}
\usepackage{xurl}
\usepackage[hidelinks]{hyperref}
\newcommand{\sref}[2]{\hyperlink{source:#1:#2}{[S#2]}}
\setlength{\emergencystretch}{3em}
\begin{document}
% problemId:30007758
\section*{Tax incidence anomalies}\label{30007758}
\subsection{Definitions and mathematical statement}
Fix comparable commodity-tax reforms in a specified market. Let \(p_h(\delta;z)\) be the expected tax-inclusive price \(h\geq0\) periods after an unanticipated tax change \(\delta\), conditional on the same pre-policy market state \(z\), relative to a counterfactual without reform. Define directional pass-through, when limits exist, by
\[\rho_+(h,z)=\lim_{\delta\downarrow0}\frac{p_h(\delta;z)-p_h(0;z)}{\delta},\qquad \rho_-(h,z)=\lim_{\delta\uparrow0}\frac{p_h(\delta;z)-p_h(0;z)}{\delta}.\]
The target is the asymmetry \(A(h,z)=\rho_+(h,z)-\rho_-(h,z)\), its persistence and its mechanism decomposition. Construct and discriminate between models with pricing adjustment costs, strategic complementarities, salience, tax evasion and demand changes, using prespecified additional measurements or interventions. Specify the tax unit, market boundaries, outcome aggregation and comparable policy durations; permanent increases and temporary cuts are different treatments. Identify the directional effects under explicit assignment and anticipation restrictions, and test out-of-sample mechanism implications. Show when a proposed mechanism generates nonzero asymmetry holding initial states fixed. Unequal raw before–after responses alone do not establish non-differentiable incidence or isolate a causal mechanism. This formulation focuses on commodity prices; analogous wage targets require a separate labor-market specification.

\par\textbf{Formulation and scope.} Catalogue-authored scoped formalization. Population, instruments, welfare objective and identification restrictions must be instantiated before claiming a resolution; source authors are not credited with these added assumptions.

\par\textbf{Open-question provenance.} An explicit published open-question statement was not verified. Sources below are supporting literature and must not be treated as a first listing.

\par\textbf{Historical priority.} An original explicit open-question statement was not verified; the cited supporting paper must not be treated as its first listing.

\subsection{Short English statement}
In the explicitly specified setting: Why do tax increases and cuts pass through differently to prices and wages, and when do those asymmetries persist? The mathematical statement above fixes the target, assumptions and scope of this catalogue formulation.

\subsection{Sources}
\begin{itemize}\raggedright
\item[\textnormal{[S1]}]\hypertarget{source:30007758:1}{} Youssef Benzarti. 2025. Tax Incidence Anomalies. Annual Review of Economics 17, 615–634; abstract.
\url{https://www.annualreviews.org/content/journals/10.1146/annurev-economics-081324-085805}.
DOI: 10.1146/annurev-economics-081324-085805.
{\small\textit{Supporting or current literature; see evidence qualification. Lists evidence challenging statutory-incidence invariance, elasticity-only incidence and symmetry. The accessed abstract does not explicitly declare the catalogue mechanism problem open. A 2024 working-paper version exists.}}
\end{itemize}
\end{document}
